% ═══════════════════════════════════════════════════════════════════════
\section{L'algorithme PageRank}
% ═══════════════════════════════════════════════════════════════════════

\begin{frame}{Contenu et popularité}
  \small Un moteur de recherche classe les pages selon deux types de
  critères :
  \begin{itemize}
    \item le contenu : présence des mots cherchés, position dans le titre
      ou le texte, proximité ;
    \item la popularité, mesurée par la structure des liens du Web.
  \end{itemize}
  PageRank (Brin et Page, 1998) relève du second type.
  \vspace{6pt}

  L'internaute aléatoire : sur une page qui contient $k$ liens, il clique
  sur l'un d'eux, chacun avec probabilité $1/k$. La proportion du temps
  qu'il passe sur une page mesure son importance.
\end{frame}

\subsection{L'internaute aléatoire}

\begin{frame}{La matrice de transition du mini-web}
  \begin{columns}[c]
    \begin{column}{0.55\textwidth}
      \small Départ sur la page $3$ (trois liens : vers $2$, $4$, $5$).
      Pour être sur la page $1$ après deux clics, il faut passer par $2$
      ou par $5$ :
      \[
        \P{X_2=1} = \tfrac12\cdot\tfrac13+\tfrac12\cdot\tfrac13 = \tfrac13.
      \]
      La page visitée est une chaîne de Markov ; $a_{ij}=1/k_j$ si la page
      $j$ (avec $k_j$ liens) pointe vers $i$, $0$ sinon :
      \[
        A = \begin{pmatrix}
          0 & \tfrac12 & 0 & 0 & \tfrac12\\[2pt]
          1 & 0 & \tfrac13 & 0 & \tfrac12\\[2pt]
          0 & \tfrac12 & 0 & \tfrac12 & 0\\[2pt]
          0 & 0 & \tfrac13 & 0 & 0\\[2pt]
          0 & 0 & \tfrac13 & \tfrac12 & 0
        \end{pmatrix}.
      \]
    \end{column}
    \begin{column}{0.43\textwidth}
      \FigBox[0.4\textheight]{%
        \begin{tikzpicture}[scale=1.15]
          \node[pageweb] (p1) at (0,2.4) {$1$};
          \node[pageweb] (p5) at (3.2,2.4) {$5$};
          \node[pageweb] (p2) at (0,0) {$2$};
          \node[pageweb] (p3) at (2.2,0) {$3$};
          \node[pageweb] (p4) at (4.4,0) {$4$};
          \draw[lien2] (p1) -- (p2);
          \draw[lien2] (p2) -- (p3);
          \draw[lien2] (p3) -- (p4);
          \draw[lien] (p3) -- (p5);
          \draw[lien] (p4) -- (p5);
          \draw[lien] (p5) -- (p1);
          \draw[lien] (p5) -- (p2);
        \end{tikzpicture}}
      \small Colonnes : liens qui sortent d'une page. Lignes : liens qui
      y entrent.
    \end{column}
  \end{columns}
\end{frame}

\begin{frame}[fragile]{Programme 4.5 --- Le PageRank du mini-web}
  \begin{columns}[T]
    \begin{column}{0.62\textwidth}
\begin{lstlisting}
def matrice_transition(liens, N):
    """liens[j] : pages vers lesquelles pointe j."""
    A = np.zeros((N, N))
    for j, sorties in liens.items():
        for i in sorties:
            A[i - 1, j - 1] = 1 / len(sorties)
    return A

liens = {1: [2], 2: [1, 3], 3: [2, 4, 5],
         4: [3, 5], 5: [1, 2]}
A = matrice_transition(liens, 5)
valeurs, vecteurs = np.linalg.eig(A)
k = np.argmin(abs(valeurs - 1))
q = np.real(vecteurs[:, k])
q = q / q.sum()
print("fois 55 :", 55 * q)  # [13. 20. 12. 4. 6.]
\end{lstlisting}
    \end{column}
    \begin{column}{0.36\textwidth}
      \small
      \begin{itemize}
        \item \texttt{eig} peut renvoyer des complexes (matrice non
          symétrique) : \texttt{np.real}.
        \item On ne sait pas où se trouve la valeur propre $1$ dans la
          liste : on cherche la plus proche.
        \item Avec $\vec p_0=(0,0,1,0,0)$ :\\
          $\vec p_1=\big(0,\tfrac13,0,\tfrac13,\tfrac13\big)$,\\
          $\vec p_2=\big(\tfrac13,\tfrac16,\tfrac13,0,\tfrac16\big)$.
      \end{itemize}
    \end{column}
  \end{columns}
\end{frame}

\begin{frame}{Le PageRank du mini-web}
  \small
  \[
    \vec q = \tfrac1{55}\,(13, 20, 12, 4, 6).
  \]
  \begin{itemize}
    \item $A$ est régulière (coefficients de $A^5$ tous $>0$) : le
      théorème 4.6 s'applique, quel que soit le départ.
    \item Classement : page $2$, page $1$, page $3$, page $5$, page $4$.
  \end{itemize}
  La deuxième ligne de $\vec q=A\vec q$ :
  \[
    q_2 = 1\cdot q_1+\tfrac13\,q_3+\tfrac12\,q_5
    = \tfrac{13}{55}+\tfrac13\cdot\tfrac{12}{55}+\tfrac12\cdot\tfrac6{55}
    = \tfrac{20}{55}.
  \]
  \begin{itemize}
    \item L'importance d'une page est la somme des « votes » des pages qui
      pointent vers elle ; chaque page partage son PageRank également
      entre ses liens sortants.
    \item La page $1$ (deux liens entrants) devance la page $3$ (deux
      aussi) : l'un de ses liens vient de la page $2$, la plus importante.
    \item Circulaire en apparence ; l'algèbre linéaire en fait la solution
      d'un système.
  \end{itemize}
\end{frame}

\begin{frame}{Note historique 4.3 --- Le gogol et un vecteur propre à 25 milliards}
  \small
  \begin{itemize}
    \item Le projet de Page et Brin sur les liens du Web donne naissance à
      Google (article de 1998).
    \item Google vient de googol (gogol en français) : le nombre
      $10^{100}$. Selon Edward Kasner (Mathematics and the Imagination,
      1940), le mot aurait été inventé par son neveu de neuf ans, Milton
      Sirotta.
    \item Kurt Bryan et Tanya Leise ont popularisé l'analyse mathématique
      de l'algorithme sous le titre « The \$25,000,000,000 eigenvector ».
      Tout repose sur les notions de ce chapitre.
  \end{itemize}
\end{frame}

\subsection{Les pages sans liens sortants}

\begin{frame}{Figure 4.5 --- Le mini-web 2}
  \begin{columns}[c]
    \begin{column}{0.4\textwidth}
      \FigBox[0.55\textheight]{%
        \begin{tikzpicture}[scale=1.05]
          \node[pageweb] (p1) at (0,2.4) {$1$};
          \node[pageweb] (p5) at (3.2,2.4) {$5$};
          \node[pageweb] (p2) at (0,0) {$2$};
          \node[pageweb] (p3) at (2.2,0) {$3$};
          \node[pageweb] (p4) at (4.4,0) {$4$};
          \node[pageweb, fill=firebrick3!15, draw=firebrick3] (p6) at (3.2,4.4) {$6$};
          \draw[lien2] (p1) -- (p2);
          \draw[lien2] (p2) -- (p3);
          \draw[lien2] (p3) -- (p4);
          \draw[lien] (p3) -- (p5);
          \draw[lien] (p4) -- (p5);
          \draw[lien] (p5) -- (p1);
          \draw[lien] (p5) -- (p2);
          \draw[lien] (p5) -- (p6);
        \end{tikzpicture}}
    \end{column}
    \begin{column}{0.58\textwidth}
      \small PDF, images, pages en construction : la page $6$ n'a aucun
      lien. La sixième colonne de $A$ est nulle : $A$ n'est plus
      stochastique, la probabilité « s'échappe ».
      \vspace{4pt}

      Correction : depuis une page sans lien, l'internaute choisit une
      page au hasard parmi les $N$ (probabilité $1/N$ chacune). Avec
      $\vec E$, l'indicatrice des pages sans lien
      ($\vec E=(0,0,0,0,0,1)$ ici) :
      \[
        B = A+\frac1N\,\vec e\,\vec E^{\mathsf T}. \tag{4.2}
      \]
    \end{column}
  \end{columns}
\end{frame}

\begin{frame}{Le PageRank du mini-web 2}
  \small
  \[
    B = \begin{pmatrix}
      0 & \tfrac12 & 0 & 0 & \tfrac13 & \tfrac16\\[2pt]
      1 & 0 & \tfrac13 & 0 & \tfrac13 & \tfrac16\\[2pt]
      0 & \tfrac12 & 0 & \tfrac12 & 0 & \tfrac16\\[2pt]
      0 & 0 & \tfrac13 & 0 & 0 & \tfrac16\\[2pt]
      0 & 0 & \tfrac13 & \tfrac12 & 0 & \tfrac16\\[2pt]
      0 & 0 & 0 & 0 & \tfrac13 & \tfrac16
    \end{pmatrix}
    \qquad \text{stochastique et régulière.}
  \]
  \[
    \vec q = \tfrac1{127}\,(27, 42, 27, 10, 15, 6)
    \approx (\num{0.213}\,;\ \num{0.331}\,;\ \num{0.213}\,;\
    \num{0.079}\,;\ \num{0.118}\,;\ \num{0.047}).
  \]
  Classement : page $2$, pages $1$ et $3$ à égalité, page $5$, page $4$,
  page $6$. Peu de changements : la page $1$ rejoint la page $3$, en
  partie parce que la page $5$ partage maintenant son vote entre trois
  liens.
\end{frame}

\subsection{Les régions piégeantes et la téléportation}

\begin{frame}{Figure 4.6 --- Le mini-web 3 : une région piégeante}
  \begin{columns}[c]
    \begin{column}{0.4\textwidth}
      \FigBox[0.55\textheight]{%
        \begin{tikzpicture}[scale=1.05]
          \node[pageweb] (p1) at (0,2.4) {$1$};
          \node[pageweb] (p5) at (3.2,2.4) {$5$};
          \node[pageweb] (p2) at (0,0) {$2$};
          \node[pageweb] (p3) at (2.2,0) {$3$};
          \node[pageweb] (p4) at (4.4,0) {$4$};
          \node[pageweb, fill=firebrick3!15, draw=firebrick3] (p6) at (3.2,4.4) {$6$};
          \node[pageweb, fill=firebrick3!15, draw=firebrick3] (p7) at (0.8,4.4) {$7$};
          \draw[lien2] (p1) -- (p2);
          \draw[lien2] (p2) -- (p3);
          \draw[lien2] (p3) -- (p4);
          \draw[lien] (p3) -- (p5);
          \draw[lien] (p4) -- (p5);
          \draw[lien] (p5) -- (p1);
          \draw[lien] (p5) -- (p2);
          \draw[lien] (p5) -- (p6);
          \draw[lien2] (p6) -- (p7);
        \end{tikzpicture}}
    \end{column}
    \begin{column}{0.58\textwidth}
      \small Une page $7$, et des liens $6\leftrightarrow7$. On entre dans
      la région $\{6,7\}$ par la page $5$, mais on ne peut plus en sortir.
      \begin{itemize}
        \item $A$ est stochastique mais pas régulière.
        \item Unique équilibre
          $\vec q=\big(0,0,0,0,0,\tfrac12,\tfrac12\big)$ : les pages $1$ à
          $5$ ont un PageRank nul. Inutilisable.
        \item La loi ne converge même pas (aller-retour entre $6$ et $7$,
          comme l'exemple 4.5) ; $-1$ est valeur propre de $A$.
        \item Sur le Web réel, de telles régions se forment spontanément,
          ou délibérément pour accaparer le PageRank.
      \end{itemize}
    \end{column}
  \end{columns}
\end{frame}

\begin{frame}{La téléportation}
  \small À chaque étape, l'internaute suit un lien avec probabilité
  $\beta$, et avec probabilité $1-\beta$ saute vers une page choisie au
  hasard parmi les $N$ (il tape une nouvelle adresse). $\beta\in\intoo01$
  est le facteur d'amortissement ; Brin et Page utilisaient
  $\beta=\num{0.85}$.
  \[
    G = \beta B+(1-\beta)\,U,\qquad U = \frac1N\,\vec e\,\vec e^{\mathsf T}
    \quad\text{(tous les coefficients valent $1/N$).} \tag{4.3}
  \]
  C'est le « tirage au sort » de l'exemple 4.3, avec la loi uniforme.
  \begin{Enonce}{Proposition 4.7 --- Propriétés de la matrice $G$}
    \small Si $0<\beta<1$, $G$ est stochastique à coefficients
    strictement positifs ; elle est donc régulière, et a un unique
    équilibre $\vec q$, à composantes strictement positives.
  \end{Enonce}
  \begin{AuTableau}
    \footnotesize Somme d'une colonne : $\beta\cdot1+(1-\beta)\cdot1=1$ ;
    chaque coefficient est au moins $(1-\beta)/N>0$ ; $k=1$.
  \end{AuTableau}
\end{frame}

\begin{frame}[fragile]{Programme 4.7 --- La matrice Google}
  \begin{columns}[T]
    \begin{column}{0.6\textwidth}
\begin{lstlisting}
def matrice_google(liens, N, beta=0.85):
    A = matrice_transition(liens, N)
    E = (A.sum(axis=0) == 0).astype(float)
    e = np.ones(N)
    B = A + np.outer(e, E) / N    # pages sans sortie
    U = np.outer(e, e) / N        # téléportation
    return beta * B + (1 - beta) * U

def equilibre(M):
    valeurs, vecteurs = np.linalg.eig(M)
    k = np.argmin(abs(valeurs - 1))
    q = np.real(vecteurs[:, k])
    return q / q.sum()

liens3 = {1: [2], 2: [1, 3], 3: [2, 4, 5],
          4: [3, 5], 5: [1, 2, 6],
          6: [7], 7: [6]}
G3 = matrice_google(liens3, 7)
print(equilibre(G3))
\end{lstlisting}
    \end{column}
    \begin{column}{0.38\textwidth}
      \small
      \begin{itemize}
        \item \texttt{np.outer(u, v)} : le produit
          $\vec u\,\vec v^{\mathsf T}$.
        \item \texttt{beta=1} : seulement la correction des pages sans
          lien (mini-web 2).
        \item Avec $\beta=\num{0.85}$, classement : $6$, $7$, $2$, $3$,
          $1$, $5$, $4$.
        \item La région piégeante reste en tête, sans monopoliser le
          classement.
      \end{itemize}
    \end{column}
  \end{columns}
\end{frame}

\begin{frame}{Tableau 4.3 --- Le choix de $\beta$}
  \small PageRank du mini-web 3 ; dernière colonne : la part de la région
  piégeante.
  \begin{center}
  \footnotesize
  \begin{tabular}{cccccccc|c}
    \toprule
    $\beta$ & $q_1$ & $q_2$ & $q_3$ & $q_4$ & $q_5$ & $q_6$ & $q_7$ & $q_6+q_7$ \\
    \midrule
    $\num{0.5}$ & $\num{0.137}$ & $\num{0.183}$ & $\num{0.141}$ & $\num{0.095}$ & $\num{0.119}$ & $\num{0.169}$ & $\num{0.156}$ & $\num{0.325}$ \\
    $\num{0.85}$ & $\num{0.122}$ & $\num{0.182}$ & $\num{0.123}$ & $\num{0.056}$ & $\num{0.080}$ & $\num{0.225}$ & $\num{0.212}$ & $\num{0.437}$ \\
    $\num{0.99}$ & $\num{0.031}$ & $\num{0.048}$ & $\num{0.031}$ & $\num{0.012}$ & $\num{0.017}$ & $\num{0.432}$ & $\num{0.429}$ & $\num{0.861}$ \\
    \bottomrule
  \end{tabular}
  \end{center}
  \small
  \begin{itemize}
    \item $\beta$ proche de $1$ : fidèle à la structure des liens, mais la
      région piégeante accapare le PageRank.
    \item $\beta$ petit : le classement se rapproche de la loi uniforme et
      perd son information.
  \end{itemize}
\end{frame}

\subsection{La matrice Google}

\begin{frame}{La matrice Google}
  \small Les deux corrections réunies :
  \[
    G = \beta B+(1-\beta)U
    = \beta A+\frac1N\,\vec e\,\big(\beta\vec E^{\mathsf T}+(1-\beta)\,\vec e^{\mathsf T}\big),
  \]
  et le PageRank est l'unique vecteur de probabilités tel que
  $G\vec q=\vec q$.
  \begin{itemize}
    \item Web réel : $N$ en milliards ; $G$ aurait plus de $10^{18}$
      coefficients, tous non nuls. Impossible à stocker.
    \item Inutile : il suffit de $A$, $\vec E$ et $\beta$.
    \item $A$ est creuse : quelques dizaines de liens par page au plus.
      On stocke les coefficients non nuls et leurs positions, c'est-à-dire
      la liste des liens.
    \item Reste à calculer $\vec q$ sans jamais former $G$.
  \end{itemize}
\end{frame}
